%https://matematicasn.blogspot.com/2016/03/inductores-ejercicios-resueltos-de.html
\documentclass[../main]{subfiles}
\begin{document}
\chapter{Inductancia Equivalente}
\img{images/10/inductores}{Inductores toroidales.}
\info{
Inductancia Equivalente}{
\subsubsection{Paralelo} 
\emph{La Inductancia equivalente es la resultante total total presente inductores de un circuito, el valor total de los inductores en paralelo es igual al recíproco de la suma de los recíprocos de las inductancias.}

\begin{equation*}
L_T= \frac{1}{\frac{1}{L_1}+\frac{1}{L_2}} \Rightarrow
L_T= \frac{1}{\frac{1}{L_1}+\frac{1}{L_2}+\dots+\frac{1}{L_n}}
\end{equation*}


\subsubsection{Serie}
\emph{La inductancia equivalente de inductores conectados en serie es igual a la suma de las coeficiente de inducción individuales.}

\begin{equation*}
L_T= L_1+L_2 \Rightarrow L_T= L_1+L_2+\dots+L_n
\end{equation*}

Equivalencias: La unidad para inductor es el Henrio $H$.

$1H = 1x10^3 m H$ (mili Henrio)

}
\actividad{Calcular la inductancia equivalente de los siguientes circuitos entre $a$ y $b$.}

$L_1=2 mH$

$L_2=4 mH$

$L_3=6 mH$

$L_4= 5 mH$

$L_5 = 6 mH$

$L_6 = 3 mH$

$L_7= 7 mH$

%\includegraphics[scale=0.3]{images/09/pizarron.jpg}
\begin{enumerate}
    \item .
    \begin{center}
        \begin{circuitikz}
         \draw (-2,0) node[label={above:a}]{} to[inductor=$L_1$,,*-] (0,0); 
         \draw (0,0) to[inductor={$L_2$}] (2,0); 
         \draw (2,0) to[inductor=$L_3$,,-*] (4,0) node[label={above:b}]{}; 
         
         \end{circuitikz}
        \end{center}
    
    \item .
        \begin{center}
        \begin{circuitikz}
         \draw (-3,0) node[label={above:a}]{} to[inductor=$L_2$,*-*](3,0) node[label={above:b}]{} (2,0)--(2,2) to[inductor=$L_1$] (-2,2)--(-2,0);
         \draw (-2,0)--(-2,-2);
         \draw (2,0)--(2,-2);
         \draw (-2,-2) to[inductor=$L_3$] (2,-2);
        \end{circuitikz}
        \end{center}
    
    \item .
    
        \begin{center}
        \begin{circuitikz}
         \draw (-3,0) node[label={above:a}]{} to[inductor=$L_1$,*-](0,0) to[inductor=$L_2$,-*] (3,0) node[label={above:b}]{} (2,0)--(2,-2) to[inductor=$L_3$] (-2,-2)--(-2,0);
   
        \end{circuitikz}
        \end{center}
    \item .
        \begin{center}
        \begin{circuitikz}
         \draw (-3,0) node[label={above:a}]{} to[inductor=$L_1$,*-*](0,3) to[inductor=$L_2$,-*] (3,0) node[label={above:b}]{}  to[inductor=$L_3$] (-3,0);
        \end{circuitikz}
        \end{center}
    \item .
        \begin{center}
        \begin{circuitikz}
         \draw (-3,0) node[label={above:a}]{} to[inductor=$L_1$,*-*](0,3) to[inductor=$L_4$,-*] (3,0) node[label={above:b}]{}  to[inductor=$L_5$] (-3,0);
        \end{circuitikz}
        \end{center}

      \item .
    
        \begin{center}
        \begin{circuitikz}
         \draw (-3,0)  to[inductor=$L_1$,*-*](0,0) node[label={above:a}]{} to[inductor=$L_2$,-*] (3,0) node[label={above:b}]{} (2,0)--(2,-2) to[inductor=$L_3$] (-2,-2)--(-2,0);
   
        \end{circuitikz}
        \end{center}
    
    \item .
        \begin{center}
        \begin{circuitikz}
         \draw (-3,0) node[label={above:a}]{} to[inductor=$L_2$,*-*](2,0)  (2,0)--(2,2) to[inductor=$L_1$] (-2,2)--(-2,0);
         \draw (-2,0)--(-2,-2);
         \draw (2,0)--(2,-2);
         \draw (-2,-2) to[inductor=$L_3$] (2,-2);
         \draw (2,0) to[inductor=$L_4$] (4,0);
         \draw (4,0) to[inductor=$L_5$,-*] (6,0) node[label={above:b}]{};
        \end{circuitikz}
        \end{center}
    
    \item .
        \begin{center}
        \begin{circuitikz}
         \draw (-3,0) node[label={above:a}]{} to[inductor=$L_1$,*-](3,0) to[inductor=$L_2$,-*] (3,-3) node[label={right:b}]{};
        \draw (-3,0)  to[inductor=$L_3$,*-](-3,-3) to[inductor=$L_4$,-*] (3,-3);
        \end{circuitikz}
        \end{center}
    
    \item .
    
      \begin{center}
        \begin{circuitikz}
        \draw (-3,0)node[label={above:a}]{} -- (-2,1) to[inductor=$L_1$,*-*](0,1)
        to[inductor=$L_2$,*-*] (2,-1)
        to[inductor=$L_3$,*-*] (0,-3)
        to[inductor=$L_4$,*-*] (-2,-3)
        -- (-2.5,-2.5)--(-2,-2)
        to[inductor=$L_5$,*-*] (-3,-1)--(-3.5,-1.5)--(-4,-1) node[label={above:b}]{};
        \draw (-2.5,-2.5)--(-3,-3)
        to[inductor=$L_6$,*-*] (-4,-2) --(-3.5,-1.5);
        
        \end{circuitikz}
        \end{center}

\item .
      \begin{center}
        \begin{circuitikz}
           \draw (-3,0) node[label={above:a}]{} to[inductor=$L_1$,*-](3,0) to[inductor=$L_4$,-*] (3,-3) node[label={right:b}]{};
        \draw (-3,0)  to[inductor=$L_2$,*-](-3,-3) to[inductor=$L_5$,-*] (3,-3);
        \draw (-3,0) to[inductor=$L_3$,-*] (3,-3);
        
        \end{circuitikz}
        \end{center}
\item .
      \begin{center}
        \begin{circuitikz}
           \draw (0,2) node[label={above:a}]{} to[inductor=$L_1$,*-*](2,4)
           to[inductor=$L_2$,-*] (4,4)
            to[inductor=$L_3$,-*] (6,2)
            to[inductor=$L_4$,-] (0,2)
            to[inductor=$L_5$,-*] (2,0)
            to[inductor=$L_6$,-*] (4,0)
            to[inductor=$L_7$,-] (6,2) node[label={above:b}]{};
        
        \end{circuitikz}
        \end{center}    
        
\item .
      \begin{center}
        \begin{circuitikz}
           \draw (0,1)--(1,0) 
           to[inductor=$L_5$,*-*] (4,0) node[label={right:b}]{}
           to[inductor=$L_3$,*-*] (4,3)node[label={right:a}]{} --(3,4)
           to[inductor=$L_2$,*-*] (0,4)
           to[inductor=$L_4$,-*] (0,1);
           \draw (0,4) to[inductor=$L_1$,*-*] (4,0); 
           
        \end{circuitikz}
        \end{center}    
\item .
              \begin{center}
        \begin{circuitikz}
           \draw (0,0) node[label={left:a}]{} to[inductor=$L_3$,*-*] (4,0) 
           to[inductor=$L_6$,*-*] (7,2)--(7,5)node[label={right:b}]{} 
           to[inductor=$L_5$,*-*] (3,5)--(3,2)
           to[inductor=$L_2$,*-*] (1.5,1)--(0,0);
           \draw (0,0) to[inductor=$L_1$,*-*] (0,4)
           to[inductor=$L_4$,*-*] (4,4)--(4,0);
           ; 
        \end{circuitikz}
        \end{center}    
\item .
              \begin{center}
        \begin{circuitikz}
           \draw (0,0) node[label={left:a}]{} to[inductor=$L_1$,*-*] (2,0);
           \draw (0,2) to[inductor=$L_2$,*-*] (2,2);
           \draw (2,0) to[inductor=$L_3$,*-*] (4,0);
           \draw (0,0) -- (0,2);
           \draw (2,0) -- (2,2);
           \draw (4,0) to[inductor=$L_4$,*-*] (4,-2);
           \draw (4,-2) to[inductor=$L_5$,*-*] (4,-4);
           \draw (4,0)--(10,0);
           \draw (6,-2)--(10,-2);
           \draw (6,0) to[inductor=$L_6$,*-*] (6,-2);
           \draw (8,0) to[inductor=$L_7$,*-*] (8,-2);
           \draw (10,0) to[inductor=$L_1$,*-*] (10,-2);
           \draw (8,-2) to[inductor=$L_1$,*-*] (8,-4);
           \draw  (4,-4) node[label={left:b}]{} -- (8,-4);
            
        \end{circuitikz}
        \end{center}    
\item .
              \begin{center}
        \begin{circuitikz}
           \draw (0,0)node[label={left:a}]{}  to[inductor=$L_6$,*-*] (2,0);
            
            \draw (2,1) -- (2,-1);
            \draw (4,1) -- (4,-1);
            
           \draw (2,1)  to[inductor=$L_1$,*-*] (4,1);
           \draw (2,-1) to[inductor=$L_2$,*-*] (4,-1);
           
           \draw (4,0) to[inductor=$L_3$,*-*] (6,0);
           \draw (6,0) to[inductor=$L_4$,*-*] (8,0);

            \draw (8,0) -- (8,-1);
           \draw (8,-1) to[inductor=$L_5$,-*] (8,-3);
           \draw (6,0)--(6,-2);
           \draw (5,-2)--(7,-2);
           \draw (5,-2) to[inductor=$L_6$,-*] (5,-4);
           \draw (7,-2) to[inductor=$L_7$,-*] (7,-4);
           \draw (8,-3)--(7,-4);
           \draw (7,-4)--(4,-4);
           \draw (4,-4) to[inductor=$L_2$,-*] (0,-4)node[label={left:b}]{};
           
        \end{circuitikz}
        \end{center}  
\item .
              \begin{center}
        \begin{circuitikz}
            \draw   (0,4) node[label={left:a}]{} --(4,4)[*-];
           \draw (0,0) node[label={left:b}]{} to[inductor=$L_1$,*-*] (4,0);
           \draw (4,4)  to[inductor=$L_2$,*-*] (4,0);
           \draw (4,4)  to[inductor=$L_3$,*-*] (8,4);
           \draw (8,4)  to[inductor=$L_4$,*-*] (8,0);
           \draw (4,0)  to[inductor=$L_5$,*-*] (8,0);
           \draw (8,0)  to[inductor=$L_3$,*-*] (12,0);
           \draw (12,0)  to[inductor=$L_2$,*-*] (12,4);
           \draw (12,4)  to[inductor=$L_7$,*-*] (8,4);
            \draw (12,0)  to[inductor=$L_6$,*-*] (8,4);
           \draw (4,0)  to[inductor=$L_6$,*-*] (8,4);
        

           
        \end{circuitikz}
        \end{center}    


        
\end{enumerate}

\end{document}